\section{Introdução}

\paragraph{}
Este documento apresenta as evidências práticas das atividades desenvolvidas no
módulo de Engenharia Social. O objetivo foi compreender como campanhas de
phishing podem ser utilizadas para obtenção de informações sensíveis em um
ambiente controlado de laboratório.

\paragraph{}
Durante as atividades foram utilizadas ferramentas especializadas para
simulação de campanhas de engenharia social, incluindo SET (Social Engineering
Toolkit), Gophish e BeEF (Browser Exploitation Framework), permitindo avaliar o
impacto de ataques baseados na interação humana.

\paragraph{}
Todas as atividades foram realizadas exclusivamente em ambiente de laboratório,
com fins educacionais e controlados, respeitando os princípios éticos e as boas
práticas de segurança da informação.
